\documentclass{article}

% Language setting
% Replace `english' with e.g. `spanish' to change the document language
\usepackage[english]{babel}

% Set page size and margins
% Replace `letterpaper' with `a4paper' for UK/EU standard size
\usepackage[letterpaper,top=2cm,bottom=2cm,left=3cm,right=3cm,marginparwidth=1.75cm]{geometry}

% Useful packages
\usepackage{amsmath}
\usepackage{amssymb}
\usepackage{graphicx}
\usepackage{inconsolata}
\usepackage{minted}
\usepackage[colorlinks=true, allcolors=blue]{hyperref}

\title{Det ultimate fysikk 1 øvingsdokumentet}
\author{Skrevet av André Hansen}

\begin{document}
\maketitle

\begin{abstract}
    Dette er et dokument med alle formler og likninger man trenger i fysikk 1. 
    Alt samlet på et sted!
\end{abstract}

\section{Konstanten}

Konstenter ja

\text{"Weins forskyvningskonstant: \underline{\hspace{5cm}}}

\text{Brukes i: \underline{\hspace{5cm}}}

\vspace{0.25cm}

\text{"Boltzman konstant: \underline{\hspace{5cm}}}

\text{Brukes i: \underline{\hspace{5cm}}}

\vspace{0.25cm}

\text{"Stefan-Boltsmann-konstant: \underline{\hspace{5cm}}}

\text{Brukes i: \underline{\hspace{5cm}}}

\vspace{0.25cm}

\text{Plancks konstant : \underline{\hspace{5cm}}}

\text{Brukes i: \underline{\hspace{5cm}}}

\vspace{0.25cm}

\text{Bohrs konstant : \underline{\hspace{5cm}}}

\text{Brukes i: \underline{\hspace{5cm}}}

\vspace{0.25cm}

\text{Lysfart i vakuum : \underline{\hspace{5cm}}}

\text{Brukes i: \underline{\hspace{5cm}}}

\vspace{0.25cm}

\text{Elementærladningen : \underline{\hspace{5cm}}}

\text{Brukes i: \underline{\hspace{5cm}}}

\vspace{0.25cm}

\text{Tyngdeakselerasjonen : \underline{\hspace{5cm}}}

\text{Brukes i: \underline{\hspace{5cm}}}

\vspace{0.25cm}

\text{Atommasseenheten : \underline{\hspace{5cm}}}

\text{Brukes i: \underline{\hspace{5cm}}}

\vspace{0.25cm}

\text{Lyspår : \underline{\hspace{5cm}}}

\text{Brukes i: \underline{\hspace{5cm}}}

\section{Formler og sånt}

\begin{align*}
    \bar{v} &= \underline{\hspace{5cm}} \\
    \bar{v(t)} &= \underline{\hspace{5cm}} \\    
    \bar{\bar{a}} &= \underline{\hspace{5cm}} \\
    \text{fartslikninga} & \underline{\hspace{5cm}} \\
    \text{posisjonslikning 1} & \underline{\hspace{5cm}} \\
    \text{posisjonslikning 2} & \underline{\hspace{5cm}} \\
    \text{tidløs likning} & \underline{\hspace{5cm}} \\
    g &= \underline{\hspace{5cm}} \\
    G &= \underline{\hspace{5cm}} \\
    \vec{G} &= \underline{\hspace{5cm}} \\
    \text{Newtons 1. lov} &= \underline{\hspace{5cm}} \\
    \vec{S} &= \underline{\hspace{5cm}} \\
    \rho &= \underline{\hspace{5cm}} \\
    \vec{B} &= \underline{\hspace{5cm}} \\
    R &= \underline{\hspace{5cm}} \\
    L &= \underline{\hspace{5cm}} \\
    \text{Newtons 3.lov} &= \underline{\hspace{5cm}} \\
    \text{Newtons 2.lov} &= \underline{\hspace{5cm}} \\
    \text{Energiloven} &= \underline{\hspace{5cm}} \\
    \text{Arbeid} &= \underline{\hspace{5cm}} \\
    \text{Arbeid ulik retning} &= \underline{\hspace{5cm}} \\
    \text{Friksjonsarbeid} &= \underline{\hspace{5cm}} \\
    \text{Effekt} &= \underline{\hspace{5cm}} \\
    \text{Potensiellenergi} &= \underline{\hspace{5cm}} \\
    \text{Verkningsgrad} &= \underline{\hspace{5cm}} \\
    \text{Kinetisk energi} &= \underline{\hspace{5cm}} \\
    \text{Mekanisk energi} &= \underline{\hspace{5cm}} \\
    \text{Setningen for kinetisk energi} &= \underline{\hspace{5cm}} \\
    \text{Setningen for bevaring av mekanisk energi} &= \underline{\hspace{5cm}} \\
    \text{Setningen for mekanisk energi} &= \underline{\hspace{5cm}} \\
    \text{Setningen for mekanisk energi} &= \underline{\hspace{5cm}} \\
    \text{Bevegelsesmengde} &= \underline{\hspace{5cm}} \\
    \text{Inpuls} &= \underline{\hspace{5cm}} \\
    \text{Inpulsloven} &= \underline{\hspace{5cm}} \\
    \text{bevegelsesmengde med to gjenstander} &= \underline{\hspace{5cm}} \\
    \text{Bevaring av bevegelsesmengde} &= \underline{\hspace{5cm}} \\
    \text{Elastisk st\o t} &= \underline{\hspace{5cm}} \\
    \text{Uelastisk st\o t} &= \underline{\hspace{5cm}} \\
    \text{Fullstendig uelastisk st\o t} &= \underline{\hspace{5cm}} \\
    \text{Kinetisk energi i elastisk st\o t} &= \underline{\hspace{5cm}} \\
    \text{Trykk} &= \underline{\hspace{5cm}} \\
    \text{Termofysikk 0. lov} &= \underline{\hspace{5cm}} \\
    \text{Kinetisk energi og temperatur} &= \underline{\hspace{5cm}} \\
    \text{Spesifikk varmekapasitet} &= \underline{\hspace{5cm}} \\
    \text{Varmekapasitet} &= \underline{\hspace{5cm}} \\
    \text{F\o rste lov i termofysikk} &= \underline{\hspace{5cm}} \\
\end{align*}

\begin{align*}
    \text{Abiatisk prosess} &= \underline{\hspace{5cm}} \\
    \text{Spesifikk smeltevarme} &= \underline{\hspace{5cm}} \\
    \text{Spesifikk fordampningsvarme} &= \underline{\hspace{5cm}} \\
    \text{Andre lov i termofysikk} &= \underline{\hspace{5cm}} \\
    \text{Varmefaktoren} &= \underline{\hspace{5cm}} \\
    \text{Maksimal varmefaktor} &= \underline{\hspace{5cm}} \\
    \text{Tversb\o lger} &= \underline{\hspace{5cm}} \\
    \text{Langsb\o lger} &= \underline{\hspace{5cm}} \\
    \text{Periode} &= \underline{\hspace{5cm}} \\
    \text{Frekvens} &= \underline{\hspace{5cm}} \\
    \text{B\o lgengde} &= \underline{\hspace{5cm}} \\
    \text{B\o lgefart} &= \underline{\hspace{5cm}} \\
    \text{Utsrolingstetthet} &= \underline{\hspace{5cm}} \\
    \text{Planckurver} &= \underline{\hspace{5cm}} \\
    \text{Wiens forskyvningslov} &= \underline{\hspace{5cm}} \\
    \text{Stefan-Boltzmanns lov} &= \underline{\hspace{5cm}} \\
    \text{Lysintensitet} &= \underline{\hspace{5cm}} \\
    \text{Solkonstanten} &= \underline{\hspace{5cm}} \\
    \text{Albedo} &= \underline{\hspace{5cm}} \\
    \text{Kvantehypotesen} &= \underline{\hspace{5cm}} \\
    \text{Bohrs f\o rste postulat} &= \underline{\hspace{5cm}} \\
    \text{Bohrs andre postulat} &= \underline{\hspace{5cm}} \\
    \text{Energiniv\aa er i hydrogenatomet} &= \underline{\hspace{5cm}} \\
    \text{B\o lgelengde i en spektralinje} &= \underline{\hspace{5cm}} \\
    \text{Fluorescens} &= \underline{\hspace{5cm}} \\
    \text{Luminescens} &= \underline{\hspace{5cm}} \\
    \text{Nukleontall} &= \underline{\hspace{5cm}} \\
    \text{Bevaring Ladning} &= \underline{\hspace{5cm}} \\
    \text{Bevaring Nukleontall} &= \underline{\hspace{5cm}} \\
    \text{Bevaring Energi} &= \underline{\hspace{5cm}} \\
    \text{Bevaring Bevegelsesmengde} &= \underline{\hspace{5cm}} \\
    \text{Bevaring leptonnummer} &= \underline{\hspace{5cm}} \\
    \text{Masseenergi} &= \underline{\hspace{5cm}} \\
    \text{Proton-proton kjernereaksjon} &= \underline{\hspace{5cm}} \\
    \text{PP-steg-1} &= \underline{\hspace{5cm}} \\
    \text{PP-steg-2} &= \underline{\hspace{5cm}} \\
\end{align*}

\begin{align*}
    \text{PP-steg-3} &= \underline{\hspace{5cm}} \\
    \text{Spektralklasser} &= \underline{\hspace{5cm}} \\
    \text{Ustr\aa lt effekt fra stjerne} &= \underline{\hspace{5cm}} \\
    \text{Degeneret materie} &= \underline{\hspace{5cm}} \\
    \text{Elektrisk ladning} &= \underline{\hspace{5cm}} \\
    \text{Strøm} &= \underline{\hspace{5cm}} \\
    \text{Spenning} &= \underline{\hspace{5cm}} \\
    \text{Resistans} &= \underline{\hspace{5cm}} \\
    \text{Resistans i seriekobling} &= \underline{\hspace{5cm}} \\
    \text{Resistans i parallellkobling} &= \underline{\hspace{5cm}} \\
    \text{Kirchhoffs f\o rste lov} &= \underline{\hspace{5cm}} \\
    \text{Kirchhoffs andre lov} &= \underline{\hspace{5cm}} \\
    \text{Joules lov} &= \underline{\hspace{5cm}} \\
    \text{Effekten i en komponent} &= \underline{\hspace{5cm}} \\
    \text{Virkningsgrad} &= \underline{\hspace{5cm}} \\
\end{align*}

\bibliographystyle{alpha}
\bibliography{sample}

\end{document}
